\begin{proof}
%Fix an error parameter $\eps\in(0,1)$.
Throughout the proof we assume that the vertex set of $G$ is $V(G)=[n]=\{1,2,\dots,n\}$ with $n>0$. Recall that every update preserves this vertex set.
Thus, every vertex is permanently identified with its integer label.
%For convenience, in the complexity bounds we write $|G|\coloneqq n$.

Assume without loss of generality that $\eps\leq 1$. Set 
\[L\coloneqq \lfloor \log_2 n\rfloor + 1\qquad \textrm{and} \qquad\xi\coloneqq\left\lceil 8\ln\tfrac{1}{\eps}\right\rceil.\]
Note that $\xi\ge 1$.

%The probabilistic argument below is for an oblivious adversary: the whole update sequence is fixed independently of the random choices made by the data structure at initialization.
%Together with the fixed-universe assumption, this is exactly what allows the sampled sets $S_{i,j}$ to remain fixed throughout the execution and what makes the graph seen at any fixed query time independent of the initial sampling.

%For every $i\in\{0,1,\dots,L\}$ let $p_i\coloneqq 2^{-(i+1)}$.
%For every pair $(i,j)\in\{0,1,\dots,L\}\times\{1,2,\dots,\xi\}$, independently for every vertex $v\in V(G)$, we put $v$ into $S_{i,j}$ with probability $p_i$.


\paragraph*{Sampling sets $S_{i,j}$.} Upon initialization of the data structure, for every pair $(i,j)\in [0,L]\times [\xi]$ we construct a set $S_{i,j}\subseteq V(G)$ by including every $v\in V(G)$ in $S_{i,j}$ with probability $p_i\coloneqq 2^{-(i+1)}$.
All these random choices are independent over all vertices and all pairs $(i,j)$. After initialization, the sets $S_{i,j}$ remain fixed forever.

We store each set $S_{i,j}$ explicitly as an $n$-bit array indexed by $[n]$.
Hence, the whole family $\{S_{i,j}\colon (i,j)\in [0,L]\times [\xi]\}$ is stored once and uses $\Oh(n(L+1)\xi)$ additional space.
Note that for every fixed subset $X\subseteq V(G)$ and every fixed pair $(i,j)$, the random variable $|S_{i,j}\cap X|$ has distribution $\mathrm{Bin}(|X|,p_i)$.

\paragraph*{Weighted sums.}
For every pair $(i,j)\in [0,L]\times [\xi]$, define functions $\wei^{(0)}_{i,j},\wei^{(1)}_{i,j}\colon V(G)\to \Z$ by
\[
\wei^{(0)}_{i,j}(v) \coloneqq
\begin{cases}
1 & \text{if } v\in S_{i,j},\\
0 & \text{if } v\notin S_{i,j},
\end{cases}
\qquad
\wei^{(1)}_{i,j}(v) \coloneqq
\begin{cases}
v & \text{if } v\in S_{i,j},\\
0 & \text{if } v\notin S_{i,j}.
\end{cases}
\]
These weight functions are fixed after the sampling upon initialization.
As $V(G)=[n]\subseteq \Z$, both $\wei^{(0)}_{i,j}$ and $\wei^{(1)}_{i,j}$ indeed map $V(G)$ to~$\Z$.

Let us give some intuition behind this definition.
By setting $\wei^{(1)}_{i,j}(v)$ to $v$ on $S_{i,j}$ and to $0$ outside, we make sure that for any $X\subseteq V(G)$ with $|X\cap S_{i,j}|=1$, the sum $\sum_{v\in X} \wei^{(1)}_{i,j}(v)$ is exactly the (identifier of the) unique vertex of $X\cap S_{i,j}$. The choice of the probabilities in the definition of the sets $S_{i,j}$ makes sure for any non-empty $X$, the probability that $|X\cap S_{i,j}|=1$ for at least one pair $(i,j)\in [0,L]\times [\xi]$ is at least $1-\eps$. Therefore, if we manage to maintain the sums $\sum_{v\in X} \wei^{(1)}_{i,j}(v)$ for all $(i,j)\in [0,L]\times [\xi]$, where $X\coloneqq \{v\mid f(G,v)>0\}$, then with high probability one of these sums is actually equal to the identifier of some vertex belonging to $X$ --- so the example we seek.


% is exactly what will later allow us to decode the witness under the event that exactly one positive vertex is sampled; outside that event, the same quotient may still produce a candidate, which is then filtered by $\Ver_f(G)$ in the case \cref{it:general-f}.
%Each instance of $\WeiSum_{f,w}(G)$ receives its own fixed weighting function at initialization.
%The common storage of the bit vectors representing the sets $S_{i,j}$ is counted separately from the space $S_R$ used by the black-box data structures.

\paragraph*{Data structure.}
We now give the description of the data structure $\RetVer_{f,\eps}[G]$. For each pair $(i,j)\in [0,L]\times [\xi]$ we maintain two instances of the $\WeiSum$ data structure: $\WeiSum_{f,\wei^{(0)}_{i,j}}[G]$ and $\WeiSum_{f,\wei^{(1)}_{i,j}}[G]$.
In case \ref{it:general-f}, we additionally maintain one instance of $\Ver_f[G]$, which will be used to filter out candidates for examples.
Every update of the dynamic graph $G$ is relayed to all the maintained instances.

\paragraph*{Query algorithm.}
When a query is issued to the current $G$, we inspect all the pairs $(i,j)\in [0,L]\times [\xi]$ in lexicographic order.
For every pair $(i,j)$ we compute the values 
\[\alpha_{i,j}\coloneqq \sum_{v\in V(G)} f(G,v)\cdot \wei^{(0)}_{i,j}(v)\qquad\textrm{and}\qquad \beta_{i,j}\coloneqq \sum_{v\in V(G)} f(G,v)\cdot \wei^{(1)}_{i,j}(v)\] by querying the maintained instances of $\WeiSum$.
If $\alpha_{i,j}=0$, then we continue to the next pair.
Since $f(G,v)\ge 0$ for all $v\in V(G)$ and $\wei^{(0)}_{i,j}(v)\in\{0,1\}$, we always have $\alpha_{i,j}\ge 0$.
Consequently, whenever $\alpha_{i,j}\neq 0$, in fact $\alpha_{i,j}>0$.
\begin{itemize}
\item Case \ref{it:binary-f}: If $\alpha_{i,j}=1$, then we return the vertex $\beta_{i,j}$, because then exactly one vertex of $S_{i,j}$ contributes a nonzero value to the sum defining $\beta_{i,j}$. Otherwise, if $\alpha_{i,j}>1$, we continue.
\item Case \ref{it:general-f}: Since $\alpha_{i,j}>0$ and therefore is nonzero, we form the rational number \[q\coloneqq \frac{\beta_{i,j}}{\alpha_{i,j}}.\]
When several positive summands contribute to the sums, this quotient need not identify any unique witness; this is why the algorithm needs to perform an explicit verification step.
If $q$ is not an integer, or if $q\notin[n]$, then we continue.
Otherwise, we query $\Ver_f(G)$ on the vertex $q$.
If $\Ver_f(G)$ confirms that $f(G,q)>0$, then we return $q$; otherwise we continue.
\end{itemize}
If no pair $(i,j)$ leads to returning a vertex, then we output $\bot$:  there is no vertex $v$ with $f(G,v)>0$.

\paragraph*{No false positives.} We now argue that when the query algorithm returns a vertex $v$, then we for sure have $f(G,v)>0$.
Suppose first that we are in the case \ref{it:binary-f} and the algorithm returns a vertex while processing some pair $(i,j)\in [0,L]\times [\xi]$.
Then $\alpha_{i,j}=1$, that is, $\sum_{v\in V(G)} f(G,v)\cdot \wei^{(0)}_{i,j}(v)=1$.
Since $\wei^{(0)}_{i,j}$ is the indicator function of $S_{i,j}$, this is equivalent to $\sum_{v\in S_{i,j}} f(G,v)=1$.
Now, every summand belongs to $\{0,1\}$ due to $f$ being binary.
Therefore, there exists exactly one vertex $x\in S_{i,j}$ such that $f(G,x)=1$, and every other vertex $v\in S_{i,j}\setminus\{x\}$ satisfies $f(G,v)=0$.
Consequently, $\beta_{i,j}=\sum_{v\in S_{i,j}} f(G,v)\cdot v=x$.
We conclude that the query algorithm returns the vertex $x$, and this vertex satisfies $f(G,x)=1>0$.
Since $x\in V(G)=[n]$, the returned value is indeed a valid vertex label.

Now suppose that we are in the case \ref{it:general-f}.
Then the algorithm returns a vertex only after first checking that the candidate $q$ is an integer in $[n]$, and then asking $\Ver_f[G]$, which confirms that the returned vertex $q$ indeed satisfies $f(G,q)>0$.
Hence every returned vertex is correct.
Note that such a returned vertex $q$ need not belong to the currently processed set $S_{i,j}$: when several positive vertices contribute, $q=\beta_{i,j}/\alpha_{i,j}$ is a weighted average of their labels and may accidentally equal the label of a vertex outside $S_{i,j}$.
This is irrelevant, because the required output is any vertex with positive $f$-value and $\Ver_f[G]$ confirms $f(G,q)>0$ before returning~$q$.

Therefore, the query algorithm returns no false positives in both cases.
If the algorithm outputs $\bot$, then it makes no positive claim, so again no false positive occurs.

\paragraph*{Probability of a false negative.} We now bound the probability that the algorithm returns a false negative.
Fix an arbitrary time step of the execution and suppose that a query is made at this time step.
Let $G$ be the graph at that moment.
Because the adversary is oblivious, $G$ depends only on the fixed sequence of updates, and not on the random family $\{S_{i,j}\colon (i,j)\in [0,L]\times [\xi]\}$ sampled at the initialization.

Define \[X\coloneqq \{v\in V(G)\mid f(G,v)>0\}\qquad\textrm{and}\qquad m\coloneqq |X|.\]
Thus, $X$ is a fixed subset of $V(G)$ from the point of view of the probability space governing the initial sampling of the family $\{S_{i,j}\colon (i,j)\in [0,L]\times [\xi]\}$.
%In particular, $X$ is independent of that sampling.

If $m=0$, then $f(G,v)=0$ for every vertex $v\in V(G)$.
Therefore, for every pair $(i,j)\in [0,L]\times [\xi]$ we have $\alpha_{i,j}=\sum_{v\in V(G)} f(G,v)\cdot \wei^{(0)}_{i,j}(v)=0$, so the algorithm returns the answer $\bot$, and this answer is correct.
Hence, only the case $m\ge 1$ requires further analysis.

%\michalin{Marker}

Set \[i^\star\coloneqq \lfloor \log_2 m\rfloor + 1\]
Since $1\le m\le n$, we have $1=\lfloor \log_2 1\rfloor + 1 \le i^\star \le \lfloor \log_2 n\rfloor + 1 = L$.
Let $p\coloneqq p_{i^\star}=2^{-(i^\star+1)}$.
Exponentiating base~$2$ yields $2^{i^\star-1}\le m < 2^{i^\star}$ and therefore \[1/4 = 2^{i^\star-1}\cdot 2^{-(i^\star+1)} \le mp < 2^{i^\star}\cdot 2^{-(i^\star+1)} = 1/2.\]
In particular, $p < 1/(2m)$.

Fix any $j\in [\xi]$ and consider the random variable $Y\coloneqq |S_{i^\star,j}\cap X|$.
Since every vertex of $X$ is included in $S_{i^\star,j}$ independently with probability $p$, we have $Y\sim\mathrm{Bin}(m,p)$.
Therefore \[\Pr(Y=1)=mp(1-p)^{m-1}.\]
Since $p<1/(2m)\le 1/2<1$, Bernoulli's inequality gives $(1-p)^{m-1}\ge 1-(m-1)p$.
As $p<1/(2m)$, we have $(m-1)p < (m-1)/(2m) < 1/2$, so $1-(m-1)p > 1/2$.
Consequently, \[(1-p)^{m-1} > 1/2.\]
Combining this with $mp\ge 1/4$, we get \[\Pr(Y=1)=mp(1-p)^{m-1} > 1/8.\]
Thus, we conclude that \[\Pr\bigl(|S_{i^\star,j}\cap X|=1\bigr) > 1/8\qquad\textrm{for every fixed }j\in [\xi].\]

The sets $S_{i^\star,1},\dots,S_{i^\star,\xi}$ are sampled independently, so the events $\bigl\{|S_{i^\star,j}\cap X|=1\bigr\}$, for $j=1,\dots,\xi$, are independent.
Hence \[\Pr\Bigl(|S_{i^\star,j}\cap X|\neq 1\textrm{ for all }j\in [\xi]\Bigr) \le \left(1-1/8\right)^\xi \le \exp\!\left(-\xi/8\right) \le \eps,\]
where the last inequality holds by the definition of $\xi$.

It remains to prove that if $|S_{i^\star,j}\cap X|=1$ for some $j\in[\xi]$, then the query algorithm returns a correct vertex no later than when it processes the pair $(i^\star,j)$.
Let $x$ be the unique element of $S_{i^\star,j}\cap X$.
Then $f(G,x)>0$ by the definition of $X$.
If the algorithm has already returned a vertex for an earlier pair, then that earlier returned vertex is correct by the no-false-positives property.
So suppose that the pair $(i^\star,j)$ is actually processed.

\begin{itemize}
\item Case \ref{it:binary-f}: Every vertex $v\in S_{i^\star,j}\setminus\{x\}$ satisfies $v\notin X$, hence $f(G,v)=0$, whereas $f(G,x)=1$ because $f$ is binary.
Therefore \[\alpha_{i^\star,j}=\sum_{v\in S_{i^\star,j}} f(G,v)=1\qquad \textrm{and}\qquad \beta_{i^\star,j}=\sum_{v\in S_{i^\star,j}} f(G,v)\cdot v=x.\]
Thus the algorithm returns a correct vertex $x$ satisfying $f(G,x)>0$.
\item Case \ref{it:general-f}: Again, every vertex $v\in S_{i^\star,j}\setminus\{x\}$ satisfies $v\notin X$, hence $f(G,v)=0$.
Therefore,
\[\alpha_{i^\star,j}=\sum_{v\in S_{i^\star,j}} f(G,v)=f(G,x)>0\qquad\textrm{and}\qquad \beta_{i^\star,j}=\sum_{v\in S_{i^\star,j}} f(G,v)\cdot v=f(G,x)\cdot x,\]
so $q=\beta_{i^\star,j}/\alpha_{i^\star,j}=x$.
Since $x\in[n]$, the algorithm queries $\Ver_f[G]$ on $x$, receives confirmation that $f(G,x)>0$, and returns $x$.
\end{itemize}

We conclude that a false negative can occur only if $|S_{i^\star,j}\cap X|\neq 1$ for every $j\in\{1,\dots,\xi\}$, and this happens with probability at most $\eps$.
Therefore, for every fixed query asked at any fixed time step, the probability that the data structure outputs $\bot$ despite $X\neq\emptyset$ is at most $\eps$.
This is exactly the false-negative guarantee claimed in the lemma statement.

\paragraph*{Complexity.} Upon initialization of the data structure, we have to sample and store all the sets $S_{i,j}$ for $(i,j)\in [0,L]\times [\xi]$; this takes $\Oh(n\log n\log \tfrac{1}{\eps})$ time. Also, we have to initialize all the $2(L+1)\xi=\Oh(\log n\log \tfrac{1}{\eps})$ maintained instances of $\WeiSum$ and, in the case \ref{it:general-f}, also the maintained instance of $\Ver$; this takes $\Oh(I\log n\log \tfrac{1}{\eps})$ time.
Therefore, the initialization time is $\Oh((n+I)\log n\log \tfrac{1}{\eps})$, as promised. Note also that the stored sets $S_{i,j}$ take $\Oh(n\log n\log \tfrac{1}{\eps})$ space, the maintained instances of $\WeiSum$ and $\Ver$ take $\Oh(M\log n\log \tfrac{1}{\eps})$ space, and no more space is used by the data structure throughout its lifetime. Therefore, the space complexity is $\Oh((M+n)\log n\log \tfrac{1}{\eps})$.

We are left with analyzing the amortized update and query complexity.
Note that every update to the maintained graph $G$ is relayed to all the $2(L+1)\xi=\Oh(\log n\log \tfrac{1}{\eps})$ maintained instances of $\WeiSum$ and, in the case \ref{it:general-f}, also to the single maintained instance of $\Ver$.
Therefore, the amortized cost of executing an update to $G$ is $\Oh(T\log n\log \tfrac{1}{\eps})$.
For every processed pair it performs two queries to maintained instances of $\WeiSum$ and, in the case \ref{it:general-f}, at most one query to $\Ver_f[G]$.
As for the amortized query complexity, the query algorithm inspects at most $(L+1)\xi=\Oh(\log n\log \tfrac{1}{\eps})$ pairs $(i,j)$, and for each inspected pair it issues two queries to the maintained instances of $\WeiSum$ and, in the case \ref{it:general-f}, at most one query to the maintained instance of $\Ver$.
Besides these calls, the algorithms performs only basic control logic of the loop and the arithmetic tests based on $\alpha_{i,j}$ and $\beta_{i,j}$, which amount to $\Oh(1)$ time per processed pair. Therefore, the amortized query time is $\Oh(T\log n\log \tfrac{1}{\eps})$ as well.
\end{proof}


%The one-time sampling and explicit storage of the family $\{S_{i,j}\}$ takes $O(n(L+1)\xi)=O\!\left(|G|\cdot (1+\log |G|)\cdot \left(1+\log(1/\eps)\right)\right)$ time and space.
%Initializing the maintained black-box structures contributes $O\!\bigl(2(L+1)\xi\cdot I_R + I_V\bigr)=O\!\bigl((L+1)\xi\cdot I_R + I_V\bigr)$ time.
%Hence the total initialization time is $O\!\bigl((L+1)\xi\cdot I_R + I_V + n(L+1)\xi\bigr)=O\!\left((I_R+I_V+|G|)\cdot (1+\log |G|)\cdot \left(1+\log(1/\eps)\right)\right)$.
%The term $|G|$ here comes solely from the explicit sampling and storage of the family $\{S_{i,j}\}$.

%For the amortized update complexity, observe that every outer update is relayed to all $2(L+1)\xi$ maintained instances of $\WeiSum$ and, in the case \cref{it:general-f}, also to the single maintained instance of $\Ver_f(G)$.
%Therefore the amortized cost of one outer update is $O\!\bigl((L+1)\xi\cdot T_R + T_V\bigr)=O\!\bigl((T_R+T_V)(L+1)\xi\bigr)=O\!\left((T_R+T_V)\cdot (1+\log |G|)\cdot \left(1+\log(1/\eps)\right)\right)$, because $(L+1)\xi\ge 1$.

%For the amortized query complexity, one outer query inspects at most $(L+1)\xi$ pairs $(i,j)$.
%For every processed pair it performs two queries to maintained instances of $\WeiSum$ and, in the case \cref{it:general-f}, at most one query to $\Ver_f(G)$.
%Hence the total time spent on black-box queries during one outer query is $O\!\bigl((Q_R+Q_V)(L+1)\xi\bigr)$.
%Besides these calls, the reduction performs only the control logic of the loop and the arithmetic tests based on $\alpha_{i,j}$ and $\beta_{i,j}$.
%Under the standard word-RAM convention, these arithmetic tests contribute only $O(1)$ time per processed pair; otherwise their cost has to be added separately.
%Thus the amortized query time is $O\!\bigl((Q_R+Q_V)(L+1)\xi\bigr)=O\!\left((Q_R+Q_V)\cdot (1+\log |G|)\cdot \left(1+\log(1/\eps)\right)\right)$.

%For the space complexity, the maintained instances of $\WeiSum$ use $O\!\bigl(2(L+1)\xi\cdot S_R\bigr)=O\!\bigl(S_R(L+1)\xi\bigr)$ space.
%In the case \cref{it:general-f}, the maintained instance of $\Ver_f(G)$ uses $O(S_V)$ space.
%The explicitly stored family $\{S_{i,j}\}$ contributes additional space $O(n(L+1)\xi)$.
%Hence the whole construction uses $O\!\bigl(S_R(L+1)\xi + S_V + n(L+1)\xi\bigr)=O\!\left((S_R+S_V+|G|)\cdot (1+\log |G|)\cdot \left(1+\log(1/\eps)\right)\right)$ space.
%The term $|G|$ here comes solely from the explicit storage of the sampled family $\{S_{i,j}\}$.

%In the case \cref{it:binary-f} we adopt the convention $I_V=T_V=Q_V=S_V=0$, exactly as in the statement.
%This proves the correctness of the construction and the quantitative bounds derived in this proof.

We remark that in the proof of \cref{lem:restore-by-sampling}, it is possible to replace explicit sampling of the sets $\{S_{i,j}\colon i\in [0,L]\times [\xi]\}$ by sampling $\xi=\Theta(\log n)$ pairwise-independent hash functions $h_1,\ldots,h_\xi\colon V(G)\to [p]$, where $p$ is a prime of magnitude $\Theta(n)$. Notably, storing each such hash function $h_i$ requires only $\Oh(\log n)$ bits, instead of the $\Oh(n\log n)$ bits needed to store the sets $S_{i,1},\ldots,S_{i,\xi}$. So at the cost of a somewhat more complicated probabilistic analysis of the query algorithm, one can reduce the contribution of the sets $S_{i,j}$ to the initialization time and the space complexity from $\Oh(n\log n\log \tfrac{1}{\eps})$ to $\Oh(\log n\log \tfrac{1}{\eps})$. However, regardless of this improvement, we still need to initialize and store $\Oh(\log n\log \tfrac{1}{\eps})$ instances of $\WeiSum$, which contributes terms $\Oh(I\log n\log \tfrac{1}{\eps})$ and $\Oh(M\log n\log \tfrac{1}{\eps})$ to the initialization time and the space complexity, respectively. As we expect $I,M\geq n$ from the implementation of $\WeiSum$, these contributions anyway dominate the potential savings and the improvement has no practical impact on the guarantees asserted in the statement of  \cref{lem:restore-by-sampling}. Therefore, we omit the details.


%We also provide in Appendix~\ref{app:restore-by-sampling-hash} a more compact and elegant alternative proof based on pairwise-independent hashing, at the cost of a technically more involved argument.
